
\documentclass[a4paper,11pt]{article}
\usepackage{fontspec}
\setmainfont{Noto Sans}
\usepackage[french]{babel}
\usepackage{microtype}
\usepackage[top=9mm,bottom=8mm,left=11mm,right=11mm]{geometry}
\usepackage{xcolor}
\usepackage{tikz}
\usetikzlibrary{arrows.meta,calc,positioning}
\usepackage[most]{tcolorbox}
\usepackage{ulem}
\pagestyle{empty}

\definecolor{Blue}{HTML}{356FA8}
\definecolor{BlueDark}{HTML}{244D78}
\definecolor{BlueLight}{HTML}{EAF2F9}
\definecolor{BlueBorder}{HTML}{B9D3E3}
\definecolor{Violet}{HTML}{7A5AA6}
\definecolor{VioletDark}{HTML}{563A7A}
\definecolor{VioletLight}{HTML}{F2ECF8}
\definecolor{VioletBorder}{HTML}{D3C4E2}
\definecolor{Neutral}{HTML}{6C7378}
\definecolor{NeutralDark}{HTML}{42494D}
\definecolor{NeutralLight}{HTML}{F0F2F3}
\definecolor{NeutralBorder}{HTML}{C9CED1}
\definecolor{Ink}{HTML}{172A32}
\definecolor{Grey}{HTML}{64757D}
\definecolor{Red}{HTML}{D52F35}
\definecolor{RedLight}{HTML}{FCEBED}

\tcbset{
 enhanced,boxrule=.6pt,arc=2.8mm,
 left=3.7mm,right=3.7mm,top=1.75mm,bottom=1.75mm,
 boxsep=1mm,coltext=Ink
}
\newtcolorbox{bluebox}[1][]{colback=BlueLight,colframe=BlueBorder,#1}
\newtcolorbox{violetbox}[1][]{colback=VioletLight,colframe=VioletBorder,#1}
\newtcolorbox{neutralbox}[1][]{colback=NeutralLight,colframe=NeutralBorder,#1}
\newtcolorbox{warnbox}[1][]{colback=RedLight,colframe=Red,
 left=2.7mm,right=2.7mm,top=1.25mm,bottom=1.25mm,#1}

\begin{document}
\begin{tikzpicture}[remember picture,overlay]
\draw[draw=BlueBorder,line width=.9pt,rounded corners=1.7mm]
([xshift=1mm,yshift=-1mm]current page.north west)--([xshift=-1mm,yshift=-1mm]current page.north east)--
([xshift=-1mm,yshift=1mm]current page.south east)--([xshift=1mm,yshift=1mm]current page.south west)--cycle;
\fill[Blue]([xshift=1mm,yshift=-1mm]current page.north west) rectangle ([xshift=5.2mm,yshift=1mm]current page.south west);
\end{tikzpicture}
\vspace*{1mm}\hspace*{6mm}
{\color{BlueDark}\fontsize{25}{27}\selectfont\bfseries IMPERFÈIT DEL SUBJONTIU}
\vspace{0.5mm}\hspace*{6mm}{\color{Grey}\large L'imparfait du subjonctif}
\vspace{2.5mm}\hspace*{6mm}\begin{minipage}{0.91\textwidth}
\begin{bluebox}
\textbf{\color{BlueDark}CONSTRUCCION}\\[1mm]\centering
{\Large \textbf{RADICAL + TERMINASONS DE L'IMPERFÈIT DEL SUBJONTIU}}\\[1mm]
\small Exemple : \textbf{CANTAR} — \textbf{cantèsse, cantèsses, cantèsse, cantèsson...}
\end{bluebox}
\vspace{2.5mm}
\begin{center}\begin{tikzpicture}[>=Latex]
\draw[Ink!65,line width=.8pt](0,0)--(15.1,0);
\node[below=2mm]at(2,0){\small PASSAT};\node[below=2mm]at(7.5,0){\small ARA};\node[below=2mm]at(13.5,0){\small FUTUR};
\draw[Blue,line width=2.5pt](2.3,0)--(6.1,0);
\draw[BlueDark,line width=1.1pt,->](7.5,1)--(7.5,.18);
\node[above=3.8mm]at(7.5,1){\bfseries\color{BlueDark}ARA};
\node[above=3mm,align=center,text width=6.2cm]at(4.2,0){\small\bfseries valor subjontiva situada dins un contèxt passat};
\end{tikzpicture}\end{center}
\vspace{1.5mm}
\begin{bluebox}\textbf{\color{BlueDark}A QUÉ SERVÍS ?}\\[1mm]
\begin{itemize}\setlength{\itemsep}{.4mm}
\item exprimir volontat, dobte, necessitat o possibilitat dins un \textbf{contèxt passat};
\item aparéisser après una proposicion principala al passat;
\item es una forma de registre puslèu \textbf{literari o soutengut} uèi.
\end{itemize}\end{bluebox}
\vspace{2mm}
\begin{bluebox}\textbf{\color{BlueDark}EXEMPLES}\\[1mm]
\textbf{Voliá que venguèsses amb nosautres.}\\
\textbf{Calguèt que partiguèssem lèu.}\\
\textbf{Dubiava que foguèsses aquí.}\\
\textbf{Volián que trabalhèsson amassa.}
\end{bluebox}
\vspace{2mm}
\begin{minipage}{.485\textwidth}\begin{bluebox}[height=3.15cm]
\textbf{\color{BlueDark}PRESENT / IMPERFÈIT}\\[1mm]
\textit{Vòli que vengas.}\\
\textit{Voliá que venguèsses.}\\[1mm]
Lo temps del subjontiu s'alinha sovent amb lo temps de la proposicion principala.
\end{bluebox}\end{minipage}\hfill
\begin{minipage}{.485\textwidth}\begin{warnbox}[height=3.15cm]
\textbf{\color{Red}ATENCION}\\[1mm]
Aquesta forma es pas sonque « lo passat del subjontiu » : sa distribucion e son usatge dependon del contèxt e del registre.
\end{warnbox}\end{minipage}
\vspace{2mm}
\begin{bluebox}\textbf{\color{BlueDark}A RETÉNER}\\[1mm]
L'imperfèit del subjontiu permet de presentar una volontat, un dobte, una necessitat o una possibilitat dins una situacion passada, amb una valor uèi sovent marcada coma literària o soutenguda.
\end{bluebox}
\vspace{1.5mm}\centering{\small\color{Grey}\textbf{Indicis frequents :} voliá que, calguèt que, dobtiava que, èra possible que, demandèt que...}
\vspace{1mm}{\scriptsize\color{Grey}Referéncia : occitan languedocian, graphia classica, vèrb'Òc / Lo Congrès.}
\end{minipage}\end{document}
